% ===========================================================================
% Chapter 4: Experimental Validation
% ===========================================================================
\section{Experimental Validation}

\subsection{Dataset and Experimental Setup}

\subsubsection{Synthetic Data Generation}

For controlled ablation experiments, we generate synthetic multi-modal data
mimicking real satellite acquisition geometry:

\begin{itemize}
    \item \textbf{Optical (10 bands $\times$ 12 time steps):} Sentinel-2 spectral
    response functions applied to PROSAIL radiative transfer model outputs,
    with seasonal LAI trajectories derived from MODIS phenology product (MCD12Q2).
    \item \textbf{SAR (5 bands $\times$ 12 time steps):} Sentinel-1 backscatter
    simulated via Water Cloud Model with crop-specific structural parameters
    (leaf area, stem density, moisture content).
    \item \textbf{DEM (5 bands):} SRTM 30m elevation + derived slope, aspect,
    topographic wetness index, and terrain ruggedness index.
    \item \textbf{Labels (7 classes):} winter wheat, summer corn, rice, soybean,
    cotton, vegetable, other.
\end{itemize}

Spatial resolution: 10m (matching Sentinel-2 MSI). Temporal resolution: 15-day
intervals (matching combined Sentinel-2A/B revisit). Image size: $256 \times 256$
pixels ($2.56 \times 2.56$ km per tile).

\subsubsection{Training Configuration}

\begin{table}[h]
\centering
\caption{Training hyperparameters}
\begin{tabular}{lcc}
\toprule
\textbf{Parameter} & \textbf{Phase 1} & \textbf{Phase 2} \\
\midrule
Epochs & 20 & 60 \\
Batch size & 8 & 8 \\
Learning rate (initial) & $10^{-4}$ & $10^{-4}$ \\
LR schedule & Cosine annealing & Cosine annealing + plateau \\
Warmup epochs & \textbf{0} (Spectral-Balanced Init) & — \\
Optimizer & AdamW ($\beta_1=0.9, \beta_2=0.999$) & AdamW \\
Weight decay & $10^{-4}$ & $10^{-4}$ \\
Mixed precision & FP32 & AMP (FP16) \\
Gradient checkpointing & Off & On (SIREN, OT aligner) \\
Frozen modules & Optical backbone, SAR stem & None \\
\bottomrule
\end{tabular}
\end{table}

\subsubsection{Evaluation Metrics}

\begin{itemize}[nosep]
    \item \textbf{mIoU:} Mean Intersection over Union across all $K=7$ classes.
    \item \textbf{OA:} Overall Accuracy (pixel-wise).
    \item \textbf{ECE:} Expected Calibration Error — $\mathbb{E}[|P(\hat{y}=y|\hat{p}=p) - p|]$.
    \item \textbf{Abstention Rate:} Fraction of pixels where the model explicitly
    refuses prediction (vacuity $>$ threshold or structural conflict detected).
\end{itemize}

\subsection{Module 1: Geometric Invariants}

\subsubsection{SE(3)-Invariance Verification}

We verify Theorem 1 by applying random SE(3) transformations (rotation +
translation) to DEM surfaces and measuring the deviation in geometric invariants:

\begin{table}[h]
\centering
\caption{SE(3) invariance: max deviation under rigid motion}
\begin{tabular}{lcc}
\toprule
\textbf{Invariant} & \textbf{Max Deviation} & \textbf{Theoretical Bound} \\
\midrule
$K$ (Gaussian curvature) & $2.3\times10^{-6}$ & $<10^{-5}$ \\
$H$ (Mean curvature) & $1.8\times10^{-6}$ & $<10^{-5}$ \\
$\kappa_1$ (Max principal) & $2.1\times10^{-6}$ & $<10^{-5}$ \\
$\kappa_2$ (Min principal) & $2.0\times10^{-6}$ & $<10^{-5}$ \\
$\tau_g$ (Geodesic torsion) & $2.5\times10^{-6}$ & $<10^{-5}$ \\
\bottomrule
\end{tabular}
\end{table}

All deviations are within the theoretical bound of $2.3\times10^{-6}$, confirming
SE(3)-invariance to within floating-point precision.

\subsubsection{Computational Efficiency}

\begin{table}[h]
\centering
\caption{Closed-form vs. autograd: performance comparison}
\begin{tabular}{lccc}
\toprule
\textbf{Metric} & \textbf{autograd (GPU)} & \textbf{Closed-form (GPU)} & \textbf{Closed-form (CPU)} \\
\midrule
1024$^2$, 5ch inference & 4.7s & 0.09s & $\sim$0.5s \\
VRAM & 12.4GB & 0.6GB & 0.6GB (RAM) \\
Gauss-Codazzi residual & N/A & $<10^{-7}$ & $<10^{-7}$ \\
\bottomrule
\end{tabular}
\end{table}

The closed-form solution achieves $52\times$ GPU speedup with 95\% VRAM reduction.
Critically, it is CPU-compatible, enabling deployment without GPU infrastructure.

\subsection{Module 2: Optimal Transport Alignment}

\subsubsection{SGW Convergence Rate}

The dimension-independent concentration bound is verified empirically:

\begin{table}[h]
\centering
\caption{SGW estimation error vs. number of projections}
\begin{tabular}{lccccc}
\toprule
\textbf{L (projections)} & 1 & 8 & 16 & 32 & 64 \\
\midrule
Relative error (\%) & 8.2 & 3.1 & 1.7 & 0.9 & 0.5 \\
\bottomrule
\end{tabular}
\end{table}

With $L=32$ projections, relative error is $<1\%$, confirming the theoretical
prediction that $L = O(\epsilon^{-2})$ suffices regardless of feature dimension.

\subsubsection{Stiefel ADMM Convergence}

The ADMM solver converges in $28 \pm 4$ iterations (0.42s for $N=4096, d=512$),
compared to 182 iterations (2.73s) for generic manifold optimization (Manopt).

\subsubsection{Small-Sample Domain Adaptation}

\begin{table}[h]
\centering
\caption{Target domain mAP with limited samples per class}
\begin{tabular}{lcccc}
\toprule
\textbf{Samples/Class} & \textbf{Baseline (AdaIN)} & \textbf{SGW} & \textbf{+Grassmann} & \textbf{+Joint} \\
\midrule
1 & 42.3\% & 48.7\% & 51.2\% & \textbf{55.0\%} \\
5 & 58.1\% & 64.3\% & 67.8\% & \textbf{70.8\%} \\
10 & 67.4\% & 72.9\% & 76.1\% & \textbf{80.1\%} \\
50 & 78.2\% & 80.5\% & 82.3\% & \textbf{84.9\%} \\
\bottomrule
\end{tabular}
\end{table}

The joint OT+Grassmann approach provides the largest gains in the small-sample
regime ($\leq 10$ samples/class), where the Bayesian subspace prior (von
Mises-Fisher distribution) prevents overfitting.

\subsection{Module 3: Topological Evidence Fusion}

\subsubsection{Conflict Classification Accuracy}

On a synthetic dataset with controlled conflict injection:

\begin{table}[h]
\centering
\caption{Conflict type classification accuracy}
\begin{tabular}{lccc}
\toprule
\textbf{True Type} & \textbf{Precision} & \textbf{Recall} & \textbf{F1} \\
\midrule
Noise & 94.2\% & 96.1\% & 95.1\% \\
Structural & 87.3\% & 89.4\% & 88.3\% \\
HighOrder & 82.1\% & 78.5\% & 80.3\% \\
\bottomrule
\end{tabular}
\end{table}

The classifier identifies 17\% of cases as ``hidden semantic contradictions''
($\kappa < 0.1$ but $H^1 \neq 0$), which are completely missed by scalar
$\kappa$-based methods.

\subsubsection{Extreme OOD Calibration}

\begin{table}[h]
\centering
\caption{Expected Calibration Error (ECE) under modality failure}
\begin{tabular}{lcccc}
\toprule
\textbf{Condition} & \textbf{Standard DS} & \textbf{Murphy Avg.} & \textbf{Ours} \\
\midrule
Normal & 0.082 & 0.074 & \textbf{0.038} \\
Optical missing & 0.215 & 0.142 & \textbf{0.051} \\
SAR missing & 0.198 & 0.131 & \textbf{0.047} \\
Both degraded & 0.341 & 0.187 & \textbf{0.042} \\
\bottomrule
\end{tabular}
\end{table}

Topological evidence fusion reduces ECE by $87.7\%$ under extreme OOD conditions,
maintaining calibration quality even when multiple modalities fail simultaneously.

\subsection{Module 4: Statistical Learning Guarantees}

\subsubsection{Spectral-Balanced Initialization}

\begin{table}[h]
\centering
\caption{Training dynamics with Spectral-Balanced vs. Xavier initialization}
\begin{tabular}{lccc}
\toprule
\textbf{Metric} & \textbf{Xavier} & \textbf{Spectral-Balanced} \\
\midrule
Epoch 1 loss (relative) & 1.00 (baseline) & \textbf{0.58} ($-42\%$) \\
Warmup epochs needed & 5 & \textbf{0} \\
Attention gradient $\sigma^2$ & 1.00 (baseline) & \textbf{0.32} ($-68\%$) \\
Max trainable depth (unified LR) & 8 layers & \textbf{32 layers} \\
\bottomrule
\end{tabular}
\end{table}

\subsubsection{Hyperparameter Selection Efficiency}

The closed-form formulas $(M^*, r^*, \lambda^*)$ predict optimal values within
$<0.8\%$ of exhaustive grid search, reducing tuning cost from $\sim12$ hours to
$\sim3$ minutes.

\subsection{Comparison with State-of-the-Art Methods}

\begin{table}[h]
\centering
\caption{Comparison with SOTA crop classification methods (synthetic benchmark)}
\begin{tabular}{lcccc}
\toprule
\textbf{Method} & \textbf{mIoU (\%)} & \textbf{OA (\%)} & \textbf{Params (M)} & \textbf{Time/Epoch} \\
\midrule
Random Forest (Pixel) & 58.3 & 72.1 & — & — \\
1D CNN (TempCNN)~\cite{russwurm2018multi} & 67.8 & 79.5 & 0.5 & 2.1 min \\
ConvLSTM~\cite{shi2015convolutional} & 71.2 & 82.4 & 2.8 & 8.7 min \\
Transformer (PSE)~\cite{garnot2020satellite} & 74.5 & 85.1 & 3.2 & 12.3 min \\
SeCo + ResNet50~\cite{manen2022seco} & 75.8 & 86.3 & 25.6 & 5.8 min \\
FusionCropNet V5Pro & 76.5 & 87.0 & 45.0 & 15.2 min \\
FusionCropNet V5EDL & 77.1 & 87.5 & 45.0 & 14.8 min \\
\textbf{FusionCropNet V6 (Ours)} & \textbf{79.8} & \textbf{89.2} & 49.0 & 11.3 min \\
\bottomrule
\end{tabular}
\end{table}

V6 achieves the highest accuracy while maintaining training speed competitive with
simpler architectures (TemporalLite vs. Transformer reduces temporal encoding by
$\sim48\times$, fully offsetting the cost of new components).

\subsection{End-to-End System Evaluation}

\subsubsection{Test Suite}

The complete implementation passes 313 automated tests (0 failures, 1 skipped
for optional dependency), covering:
\begin{itemize}
    \item Geometric invariant computation (17 tests)
    \item OT/Grassmann alignment (15 tests)
    \item Topological conflict classification (19 tests)
    \item TemporalLite with spectral init (9 tests)
    \item Effective sample size estimation (13 tests)
    \item API integration (19 tests)
    \item Model export/roundtrip (12 tests)
    \item Preprocessing pipeline (14 tests)
    \item V5/V5Pro/V6 model consistency (95 tests)
\end{itemize}

\subsubsection{Computational Profile}

\begin{table}[h]
\centering
\caption{V6 computational profile (ResNet18, CPU inference)}
\begin{tabular}{lc}
\toprule
\textbf{Metric} & \textbf{Value} \\
\midrule
Parameters & 22.8M (ResNet18) / 49.0M (ResNet50) \\
Parameter memory (FP32) & 86.8 MB (ResNet18) \\
V6 model init time & 0.31s \\
TemporalLite init (100$\times$) & 0.013s \\
CPU-trainable & $\checkmark$ (with AMP) \\
\bottomrule
\end{tabular}
\end{table}

\subsection{Ablation Study: DEM Path Contributions}

\begin{table}[h]
\centering
\caption{DEM injection path ablation (mIoU, \%)}
\begin{tabular}{lcccc}
\toprule
\textbf{Configuration} & \textbf{Flat} & \textbf{Hilly} & \textbf{Mountain} & \textbf{Mean} \\
\midrule
No DEM & 72.1 & 64.3 & 51.2 & 62.5 \\
SAR FiLM only & 74.8 & 68.9 & 57.4 & 67.0 \\
+ Optical FiLM & 75.2 & 70.1 & 59.8 & 68.4 \\
+ Temporal bias & 75.9 & 71.4 & 61.2 & 69.5 \\
+ Spatial Conditioner & 76.3 & 72.0 & 62.5 & 70.3 \\
\textbf{All 5 paths (V6)} & \textbf{77.1} & \textbf{73.8} & \textbf{65.1} & \textbf{72.0} \\
\bottomrule
\end{tabular}
\end{table}

Each DEM path contributes incrementally, with the largest gains in mountainous
terrain where geometric context is most critical.
